\section{Теоремы о спектральном сжатии и коллапсе персонализации}
\label{sec:Chapter3} \index{Chapter3}

% -----------------------------------------------------------------------
% 3.1 КОРРЕКТНОСТЬ ЛОКАЛЬНОЙ ДИНАМИКИ
% -----------------------------------------------------------------------

\subsection{Корректность локальной динамики}
\label{subsec:well_posedness}

Прежде чем формулировать результаты о сходимости, установим корректность
построенной модели: существование стационарных точек и применимость
линеаризации в их окрестностях.

\begin{proposition}[Существование стационарной точки]
\label{prop:fixed_point}
При выполнении условий \textup{(A1)--(A3)} (компактность пространств
$\mathcal{U}$, $\mathcal{I}$, равномерная сходимость эмпирического
функционала) отображение $\Phi = \mathcal{G} \circ \mathcal{F}$
допускает по крайней мере одну неподвижную точку
$(\bar\mu, \bar\nu) \in \mathcal{P}(\mathcal{U}) \times
\mathcal{P}(\mathcal{I})$.
\end{proposition}

\textit{Доказательство.} Применение теоремы Какутани к компактному
выпуклому пространству мер и полунепрерывному снизу оператору $\Phi$.
\hfill$\square$

Содержательно предложение~\ref{prop:fixed_point} означает, что
рекомендательная система всегда достигает некоторого режима
«самосогласованности»: существует политика и соответствующее
распределение данных, при которых переобучение не меняет политику.
Вопрос в том, является ли этот режим разнообразным или деградировавшим.

% -----------------------------------------------------------------------
% 3.2 СПЕКТРАЛЬНОЕ СЖАТИЕ ОПЕРАТОРА ПЕРЕОБУЧЕНИЯ
% -----------------------------------------------------------------------

\subsection{Спектральное сжатие оператора переобучения}
\label{subsec:spectral_contraction}

Центральный технический результат работы~--- теорема о спектральном
сжатии.

\begin{theorem}[Спектральное сжатие]
\label{thm:spectral_contraction}
Пусть выполнены условия \textup{(A1)--(A5)}. Пусть шаг градиентного
спуска $\eta$ удовлетворяет $\eta < 2/\lambda_{\max}(\mathbf{H})$,
и пусть параметр exploration равен нулю ($\alpha = 0$). Тогда для
эффективного оператора переобучения кластера $k$ выполнено:
\[
    \rho(\mathbf{A}_k) \;\leq\;
    1 - \eta\,\lambda_{\min}^+(\mathbf{H}\big|_{\mathcal{C}_k}),
\]
где $\lambda_{\min}^+$ обозначает минимальное положительное собственное
значение.
\end{theorem}

\textit{Доказательство.} Следует из явного представления
$\mathbf{A}_k \approx \mathbf{I} - \eta\,\mathbf{H}\big|_{\mathcal{C}_k}$
и оценки нормы возмущения по теореме Вейля. Подробный вывод~---
в приложении~A. \hfill$\square$

Если кластерный Хессе не вырожден (что выполнено, например, при
$L_2$-регуляризации с коэффициентом $\lambda > 0$), то
$\lambda_{\min}^+(\mathbf{H}\big|_{\mathcal{C}_k}) \geq \lambda$,
и $\rho(\mathbf{A}_k) \leq 1 - \eta\lambda < 1$. Это означает
экспоненциальное затухание ковариаций.

% -----------------------------------------------------------------------
% 3.3 ЗАТУХАНИЕ ЛОКАЛЬНОГО ЯКОБИАНА
% -----------------------------------------------------------------------

\subsection{Затухание локального Якобиана и исчезновение разнообразия}
\label{subsec:jacobian_decay}

Из теоремы~\ref{thm:spectral_contraction} немедленно следует результат
об убывании внутрикластерного разнообразия.

\begin{corollary}[Затухание разнообразия]
\label{cor:diversity_decay}
При условиях теоремы~\ref{thm:spectral_contraction} и $\alpha = 0$
внутрикластерное разнообразие удовлетворяет:
\[
    D^{(t)} \;\leq\; D^{(0)}\,
    \bigl(1 - \eta\,\lambda_{\min}^+\bigr)^{2t}
    \;\xrightarrow[t \to \infty]{}\; 0.
\]
\end{corollary}

Следствие~\ref{cor:diversity_decay} означает, что даже если на момент
$t=0$ пользователи внутри каждого кластера имеют разнообразные
предпочтения (ненулевая $\Sigma_k^{(0)}$), повторяющееся переобучение
стирает это разнообразие экспоненциально быстро.

% -----------------------------------------------------------------------
% 3.4 ТЕОРЕМА О КОЛЛАПСЕ ПЕРСОНАЛИЗАЦИИ (ЭХО-КАМЕРА)
% -----------------------------------------------------------------------

\subsection{Теорема о коллапсе персонализации}
\label{subsec:collapse_theorem}

Главный теоретический результат работы.

\begin{theorem}[Коллапс персонализации]
\label{thm:personalization_collapse}
Пусть выполнены условия теоремы~\ref{thm:spectral_contraction},
и пусть дополнительно функция отклика $\mathcal{G}$ непрерывна по
распределению рекомендаций в слабой топологии. Тогда при $\alpha = 0$
и любом начальном состоянии $(\mu^{(0)}, \nu^{(0)})$ распределение
рекомендаций сходится к $K$-модальному аттрактору:
\[
    \pi^{(t)} \;\xrightarrow[t \to \infty]{w}\;
    \pi^* \;=\; \sum_{k=1}^{K} \mu_k^*\,\delta_{i^*_k},
\]
где $\{i^*_k\}$~--- набор из $K$ <<стационарных>> объектов, по одному
на каждый кластер, а $\mu^*_k$~--- предельные веса кластеров.
\end{theorem}

\textit{Набросок доказательства.} Из следствия~\ref{cor:diversity_decay}
получаем $\Sigma_k^{(t)} \to 0$, откуда~--- слабую сходимость
внутрикластерного распределения пользовательских предпочтений к дельта-мере
в центроиде $\bar{x}_k$. Непрерывность $\mathcal{G}$ гарантирует, что
политика $\pi^{(t)}$ сходится к политике на точечных предпочтениях
$\{\bar{x}_k\}$, которая выдаёт единственный объект $i^*_k =
\arg\max_{i} f(\bar{x}_k, i;\, w^*)$ для каждого кластера. \hfill$\square$

\textbf{Интерпретация.} Теорема~\ref{thm:personalization_collapse}
формализует интуицию эхо-камеры: в пределе система перестаёт
различать пользователей внутри кластера и рекомендует одно и то же
«кластерное» содержимое. Разнообразие на уровне кластеров сохраняется
($K > 1$), но индивидуальная персонализация исчезает.

% -----------------------------------------------------------------------
% 3.5 СКОРОСТЬ СХОДИМОСТИ
% -----------------------------------------------------------------------

\subsection{Оценка скорости сходимости}
\label{subsec:convergence_rate}

\begin{proposition}[Скорость коллапса]
\label{prop:rate}
Под условиями теоремы~\ref{thm:personalization_collapse} расстояние
Васерштейна между $\pi^{(t)}$ и $\pi^*$ убывает не медленнее, чем:
\[
    W_2(\pi^{(t)},\, \pi^*) \;\leq\; C \cdot
    \rho_{\max}^t, \quad
    \rho_{\max} = \max_k \rho(\mathbf{A}_k) < 1,
\]
где константа $C$ зависит от начального состояния и геометрии кластеров.
\end{proposition}

Предложение~\ref{prop:rate} даёт конкретную характеристику скорости
деградации персонализации в терминах наблюдаемой величины~---
спектрального радиуса $\mathbf{A}_k$, который может быть оценён по
данным обучения. Это обеспечивает связь теоретической модели с
практическим мониторингом систем.
